\documentclass[conference]{IEEEtran}

\usepackage[T1]{fontenc}
\usepackage[utf8]{inputenc}
\usepackage{cite}
\usepackage{amsmath,amssymb}
\usepackage{booktabs}
\usepackage{tabularx}
\usepackage{array}
\usepackage{multirow}
\usepackage{enumitem}
\usepackage{float}
\usepackage{hyperref}
\usepackage{xcolor}
\usepackage{url}

\hypersetup{
  colorlinks=true,
  linkcolor=black,
  citecolor=black,
  urlcolor=blue
}

% ── column types ─────────────────────────────────────────────────────────────
% C: centred fixed-width; P: left-aligned fixed-width; T: tiny-font path cell
% T uses \scriptsize + \raggedright; cells should use \path{} for breakable paths
\newcolumntype{P}[1]{>{\raggedright\arraybackslash}p{#1}}
\newcolumntype{T}[1]{>{\scriptsize\raggedright\arraybackslash}p{#1}}

\title{A CAPIF-Aware AI-Native Telco Orchestration Platform for\\
Intent-to-API Planning Across CAMARA and NEF Services}

\author{
\IEEEauthorblockN{T\"urkan Do\u{g}a Durak}
\IEEEauthorblockA{
Centre Tecnol\`ogic de Telecomunicacions de Catalunya (CTTC)\\
Castelldefels, Spain\\
\texttt{turkandoga.durak@cttc.es}
}
}

\begin{document}

\maketitle

% ─────────────────────────────────────────────────────────────────────────────
\begin{abstract}
This paper presents the design, implementation, and empirical evaluation of a
mini telco orchestration platform that combines 3GPP Common API Framework
(CAPIF) integration, CAMARA-aligned service semantics, NEF-backed telecom
realization, and large language model (LLM)-assisted intent planning.
The system is implemented as a FastAPI-based control plane with a
product-style operator interface, embedded OpenAPI console, live CAPIF
invoker onboarding and service discovery, provider publication, persistent
NEF data services, and guarded orchestration execution.
The platform supports both deterministic and LLM-assisted planning modes and
validates generated API payloads through a three-layer deterministic validator
before execution.
Empirical evaluation on a 60-entry bilingual (English/Turkish) ground-truth
dataset yields: IWSR\,=\,100\% (deterministic) and IWSR\,=\,95\%
(LLM-assisted, Claude Haiku\,4.5), with SCR\,=\,SVR\,=\,100\% in both
modes and a Stability Score of 0.938 over repeated LLM runs.
The resulting system provides a practical and measurable baseline for
reliability-aware intent-to-API orchestration research in 5G and beyond-5G
service exposure environments.
\end{abstract}

\begin{IEEEkeywords}
CAPIF, CAMARA, NEF, 5G service exposure, intent-based orchestration,
LLM planning, telecom APIs, OpenCAPIF, reliability evaluation
\end{IEEEkeywords}

% =============================================================================
\section{Introduction}
% =============================================================================

Programmable network exposure is becoming a central design principle in modern
telecom systems. CAPIF provides a standardised governance and discovery plane
for exposed APIs~\cite{3gpp_capif}, CAMARA defines developer-facing service
semantics~\cite{camara2022}, and NEF-based providers bridge those interfaces
to telecom data and network capabilities~\cite{nefsim2022}.
At the same time, large language models (LLMs) introduce a new opportunity:
translating high-level natural language intents into structured API workflows
without requiring the operator to manually compose payloads and endpoint
sequences~\cite{maatouk2025,flowmind2023}.

This project implements a working telco orchestration platform at that
intersection. Rather than remaining at the conceptual architecture level,
the system integrates live OpenCAPIF onboarding and discovery, live provider
publication, live QoD and Location Retrieval execution, an embedded operator
interface, and an LLM-assisted planning layer protected by deterministic
validation. The platform is evaluated using six formally defined reliability
metrics on a bilingual intent dataset, providing empirical evidence for the
feasibility and limitations of LLM-driven API orchestration under
telecom-specific constraints.

% =============================================================================
\section{Related Work}
% =============================================================================

Intent-based networking has been explored at the NFV orchestration level by
Rafiq et al.~\cite{rafiq2020}, who demonstrated deterministic intent-to-descriptor
translation across multiple orchestrators, and surveyed by Saraiva de~Sousa
et al.~\cite{saraiva2019} as a key future direction for adaptive network
service management. Coronado et al.~\cite{coronado2022} identified intent
abstraction and AI-native closed-loop control as foundational pillars of
zero-touch 6G management.

On the exposure side, Molina~Sanchez et al.~\cite{molina2022} implemented
the CAPIF Core Function as an open-source microservice; Fragkos et al.~\cite{nefsim2022}
introduced NEFSim as an open NEF experimentation platform; and Turkovic et al.~\cite{turkovic2025}
demonstrated 3GPP-compliant CAMARA QoD and Location API integration in a real
CAM testbed, measuring call-flow latency of 230--340~ms.

LLM reliability for structured API output has been benchmarked by Liu et al.~\cite{liu2024},
who found that even GPT-4 achieves an average structured-output compliance
score of only 0.40. Maatouk et al.~\cite{maatouk2025} identified hallucination
and output inconsistency as primary risks for LLM deployment in telecom
operations. Lee et al.~\cite{telagentbench2025} introduced TelAgentBench,
a multi-faceted benchmark for LLM-based telecom agents, confirming that
multi-turn planning tasks significantly degrade pass rates.
Zeng et al.~\cite{flowmind2023} demonstrated near-perfect workflow generation
accuracy when LLMs are grounded in structured API descriptions.

To the best of our knowledge, no prior work provides empirical reliability
evaluation of LLM-based intent-to-API translation within a live
CAPIF-compliant 5G service exposure environment under telecom-specific
schema, semantic, and registry constraints.

% =============================================================================
\section{System Objectives and Contributions}
% =============================================================================

The platform was built with four primary objectives:

\begin{itemize}[leftmargin=1.1em]
  \item to unify CAPIF onboarding, discovery, and publication inside one
        operator-facing control plane,
  \item to expose CAMARA-aligned telecom capabilities through a live
        orchestration backend,
  \item to support both deterministic and LLM-assisted planning for
        natural-language requests,
  \item to ensure that execution remains guarded by deterministic validation
        and schema enforcement.
\end{itemize}

Its main contributions are:

\begin{itemize}[leftmargin=1.1em]
  \item a working CAPIF-aware telco orchestration backend,
  \item a pluggable multi-provider LLM planning layer grounded by local
        service metadata and CAPIF discovery,
  \item a three-layer deterministic validator (schema, semantic, registry),
  \item a live provider integration across QoD, Location Retrieval,
        QoS Provisioning, and QoS Profiles,
  \item six formally defined reliability metrics with a 60-entry bilingual
        ground-truth evaluation dataset,
  \item a structured feedback correction loop for bounded iterative
        repair of invalid LLM outputs.
\end{itemize}

% =============================================================================
\section{Platform Architecture}
% =============================================================================

\subsection{High-Level Components}

\begin{table}[H]
\caption{Implemented Platform Components}
\label{tab:components}
\centering
\small
\begin{tabularx}{\linewidth}{P{0.25\linewidth} X}
\toprule
\textbf{Component} & \textbf{Responsibility} \\
\midrule
Operator Interface  & Product-style landing page with embedded OpenAPI console \\
FastAPI Backend     & Route handling, orchestration control, validation, execution \\
Intent Engine       & Deterministic keyword-based intent-to-step routing \\
LLM Planner         & Two-phase structured planning under catalog and CAPIF grounding \\
Multi-Layer Validator & L1 schema, L2 semantic policy, L3 CAPIF registry checks \\
Feedback Engine     & Bounded correction loop (max 3 iterations) for invalid payloads \\
CAPIF Client        & Invoker onboarding, service discovery, provider publication \\
Provider Layer      & QoD, Location Retrieval, QoS Provisioning, QoS Profiles \\
Persistence Layer   & MongoDB-backed NEF data with automatic initialization \\
\bottomrule
\end{tabularx}
\end{table}

\subsection{Repository Structure}

\begin{table}[H]
\caption{Key Repository Paths}
\label{tab:repo}
\centering
\small
\begin{tabularx}{\linewidth}{T{0.36\linewidth} X}
\toprule
\textbf{Path} & \textbf{Role} \\
\midrule
\path{backend/main.py}            & FastAPI entry point and route registration \\
\path{backend/intent_engine.py}  & Deterministic planner and field extractor \\
\path{backend/llm_planner.py}    & Two-phase LLM planning with RAG grounding \\
\path{backend/validator.py}       & Three-layer validation (L1/L2/L3) \\
\path{backend/feedback_engine.py} & Bounded correction loop \\
\path{backend/capif/client.py}    & CAPIF onboarding, discovery, publish \\
\path{backend/evaluation_engine.py} & IWSR/SCR/SVR/HR/SS/RRF computation \\
\path{frontend/index.html}        & Operator landing page with API console \\
\path{providers/NEF/}             & NEF Monitoring Event stack + Mongo persistence \\
\path{providers/camara-QoD/}      & CAMARA Quality-on-Demand service \\
\path{providers/CamaraLocationRetrieval/} & CAMARA Location Retrieval service \\
\path{providers/NEF-QoS/}         & 3GPP AsSessionWithQoS adapter (port 8585) \\
\bottomrule
\end{tabularx}
\end{table}

% =============================================================================
\section{Backend Orchestration Design}
% =============================================================================

\subsection{Core API Surface}

\begin{table}[H]
\caption{Main Backend Endpoints}
\label{tab:endpoints}
\centering
\small
\begin{tabularx}{\linewidth}{P{0.10\linewidth} T{0.32\linewidth} X}
\toprule
\textbf{Method} & \textbf{Path} & \textbf{Purpose} \\
\midrule
GET  & \path{/health}                    & Runtime health summary \\
GET  & \path{/catalog/services}          & CAMARA-aligned capability catalog \\
POST & \path{/orchestrate}               & Intent-to-execution planning \\
GET  & \path{/capif/status}              & CAPIF connectivity check \\
GET  & \path{/capif/invoker}             & Stored invoker identity \\
POST & \path{/capif/invokers/onboard}    & Live invoker onboarding \\
GET  & \path{/capif/discover}            & Live CAPIF service discovery \\
POST & \path{/capif/publish}             & Single service publication \\
POST & \path{/capif/publish/catalog}     & Publish all catalog services \\
\bottomrule
\end{tabularx}
\end{table}

\subsection{Supported Telecom Capabilities}

The platform exposes four CAMARA-standardised orchestration targets:
Quality on Demand (QoD) for temporary QoS session creation;
QoS Profiles for profile discovery; QoS Provisioning for long-lived
QoS assignment via 3GPP AsSessionWithQoS (TS\,29.122); and Location
Retrieval for device location access through NEF-backed monitoring.

\subsection{Planning Modes}

\texttt{POST /orchestrate} supports three modes:
\textit{deterministic} (rule-based, always available);
\textit{llm-assisted} (requires configured LLM API key, fails
explicitly if planning fails); and
\textit{auto} (attempts LLM first, falls back to deterministic).
This keeps the AI layer optional yet operational for both live
demonstrations and controlled experimentation.

% =============================================================================
\section{CAPIF Integration Workflow}
% =============================================================================

\subsection{CAPIF Roles}

CAPIF provides the exposure governance plane. The platform acts as an
\textbf{Invoker} consuming exposed APIs; CAMARA providers act as
\textbf{AEFs} publishing service descriptions; and the OpenCAPIF stack
provides the \textbf{Core Function} managing discovery, authentication,
and onboarding~\cite{molina2022}.

\subsection{Invoker Onboarding Flow}

The onboarding implementation in \texttt{backend/capif/client.py}:
(1)~authenticates against the register service;
(2)~generates an RSA private key and certificate signing request (CSR)
locally;
(3)~submits the onboarding request to CAPIF Core;
(4)~receives an \texttt{apiInvokerId} and invoker certificate chain;
(5)~stores runtime credentials for subsequent discovery and authenticated
execution.
This enables the platform to operate as a real CAPIF invoker rather than
a static mock client. CAPIF thus serves three distinct functions:
\textbf{identity onboarding}, \textbf{service governance} (APIs are
visible only after publication), and \textbf{service discovery} for
orchestration-time reasoning.

\subsection{Discovery and Publication}

\texttt{GET /capif/discover} queries the CAPIF registry using the invoker's
mTLS certificate and returns live service descriptions. Provider publication
follows the standard APF sequence: onboard provider identity, submit OpenAPI
metadata, verify appearance in discovery results.

% =============================================================================
\section{CAMARA and NEF Service Realization}
% =============================================================================

CAMARA defines the API-facing service contracts (the \emph{semantic layer}),
while CAPIF remains the governance and exposure layer. Provider services from
the FRONT-research-group~\cite{turkovic2025} bridge these APIs to telecom
behavior: the QoD provider handles temporary session creation; the NEF
Monitoring Event provider resolves phone numbers to IMSI identifiers,
reads location data from MongoDB, enriches responses with polygon area
information, and returns a 3GPP-style monitoring result that the
CamaraLocationRetrieval adapter normalises into the CAMARA Location
Retrieval format~\cite{nefsim2022}.

The Location Retrieval end-to-end path executes as:
(1)~intent mapped to \texttt{location-retrieval};
(2)~provider constructs a 3GPP MonitoringEvent subscription;
(3)~provider obtains CAPIF security token;
(4)~NEF API is called;
(5)~NEF resolves identifiers and polygon data;
(6)~response normalised to CAMARA format and returned.

% =============================================================================
\section{LLM-Assisted Planning Layer}
% =============================================================================

\subsection{Two-Phase Planning Strategy}

\texttt{backend/llm\_planner.py} implements a two-phase pipeline.
\textit{Phase\,1} (service selection): the LLM receives the user intent,
a RAG-injected summary of all CAMARA services and operations, and
optionally live CAPIF discovery results; it returns a JSON object
specifying \texttt{service\_id}, \texttt{operation\_id}, \texttt{method},
\texttt{path}, and \texttt{rationale}.
\textit{Phase\,2} (payload generation): given the selected operation,
the LLM receives the full CAMARA OpenAPI schema for that operation
(parsed from YAML at runtime) and returns a complete, schema-aligned
JSON payload.
The planner proposes structured actions but never executes them directly;
all output passes through the three-layer validator before dispatch.

Three LLM providers are supported via a unified dispatcher:
Gemini (\texttt{GEMINI\_API\_KEY}), OpenAI (\texttt{OPENAI\_API\_KEY}),
and Anthropic (\texttt{ANTHROPIC\_API\_KEY}).
The full 60-entry evaluation in this paper used
\textbf{Claude Haiku\,4.5} (Anthropic, model:
\texttt{claude-haiku-4-5-20251001}, temperature\,0.1) at a total API
cost of approximately \$0.35.

\subsection{Structured Feedback Correction Loop}

When Phase\,2 output fails validation, the feedback engine constructs
a correction prompt containing the original intent, the rejected payload,
and layer-tagged error messages (e.g., \texttt{[L1] device.phoneNumber
must match E.164}), and resends it to the same LLM. This loop repeats
up to \texttt{MAX\_ITERATIONS\,=\,3}. The recovery result is recorded
in \texttt{FeedbackMetadata} and returned in the orchestration response.

\subsection{Representative Orchestration Result}

An LLM-assisted dry-run request with intent
``\textit{Boost video QoS for +905551112233 for 15~minutes}''
produced:
service\,=\,\texttt{qod},
operation\,=\,\texttt{createSession},
path\,=\,\texttt{/quality-on-demand/v1/sessions},
duration\,=\,900\,s,
profile\,=\,\texttt{QOS\_E\_STREAMING}.
The three-layer validator confirmed structural, semantic, and registry
compliance in a single pass.

% =============================================================================
\section{Reliability Evaluation Framework}
% =============================================================================

\subsection{Formal Metric Definitions}

Intent-to-API translation is formalised as a mapping
$f_\theta : \mathcal{I} \rightarrow \mathcal{W}$,
where $\mathcal{I}$ is the set of natural language intents and
$\mathcal{W}$ is the set of CAPIF-compliant API workflows generated
by LLM $f_\theta$.

\subsubsection{Intent-to-Workflow Success Rate (IWSR)}
\begin{equation}
\text{IWSR} = \frac{|\{w \in \mathcal{W} \mid w = w^*\}|}{|\mathcal{I}|}
\end{equation}
where $w^*$ is the ground-truth workflow. IWSR measures end-to-end
correctness: correct service selection AND valid schema AND valid
telecom semantics.

\subsubsection{Schema Compliance Rate (SCR)}
\begin{equation}
\text{SCR} = \frac{|\{w \in \mathcal{W} \mid w \models \text{Schema}\}|}{|\mathcal{W}|}
\end{equation}
Fraction of generated payloads passing Layer\,1 (Pydantic) validation.

\subsubsection{Semantic Validity Rate (SVR)}
\begin{equation}
\text{SVR} = \frac{|\{w \in \mathcal{W} \mid w \models \text{TelecomPolicy}\}|}{|\{w \in \mathcal{W} \mid w \models \text{Schema}\}|}
\end{equation}
Among schema-compliant payloads, fraction also passing Layer\,2
telecom policy checks (E.164 format, valid QoS profile, duration
bounds, etc.).

\subsubsection{Hallucination Rate (HR)}
\begin{equation}
\text{HR} = \frac{|\{w \in \mathcal{W} \mid w \not\subseteq \mathcal{A}_{\text{reg}}\}|}{|\mathcal{W}|}
\end{equation}
where $\mathcal{A}_{\text{reg}}$ is the set of APIs discovered via
CAPIF. Measures service-selection hallucinations (service does not
exist in the CAPIF registry).

\subsubsection{Stability Score (SS)}
Given $k$ repeated runs producing $d$ distinct output workflows:
\begin{equation}
\text{SS} = 1 - \frac{d - 1}{k}
\end{equation}
Perfect stability ($d = 1$, identical output every run) yields
$\text{SS} = 1$. Computed by running the same 60-entry dataset
5 times and measuring output variance per entry.

\subsubsection{Recovery Rate after Feedback (RRF)}
\begin{equation}
\text{RRF} = \frac{|\{w_c \mid w_i \not\models \text{Schema}
             \,\wedge\, w_c = w^*\}|}
             {|\{w_i \mid w_i \not\models \text{Schema}\}|}
\end{equation}
Among initially-invalid payloads, fraction successfully corrected
by the feedback loop to a valid, ground-truth-matching workflow.

\subsection{Evaluation Dataset}

The ground-truth dataset (\texttt{datasets/ground\_truth.json}) contains
60 manually crafted entries:
20 for QoD, 15 for Location Retrieval, 10 for QoS Profiles,
10 for QoS Provisioning, and 5 edge cases;
difficulty split: 24 easy, 20 medium, 16 hard;
language split: 32 English, 28 Turkish (50/50 bilingual coverage).
Phone numbers span four country codes (+30\,GR, +90\,TR, +34\,ES,
+1\,US) from FRONT-research-group test fixtures.

\subsection{Evaluation Results}

\begin{table}[H]
\centering
\caption{Reliability Metrics: Deterministic vs.\ LLM-Assisted (60 entries)}
\label{tab:results}
\small
\begin{tabular}{lrr}
\toprule
\textbf{Metric} & \textbf{Deterministic} & \textbf{LLM (Claude Haiku\,4.5)} \\
\midrule
IWSR  & 100.0\% & 95.0\% \\
SCR   & 100.0\% & 100.0\% \\
SVR   & 100.0\% & 100.0\% \\
HR    &   0.0\% &   5.0\% \\
SS    &   1.000 &   0.938 \\
RRF   &   N/A   &   N/A\textsuperscript{a} \\
Avg.\ latency & $<$1\,ms & $\approx$3\,s \\
\bottomrule
\end{tabular}

\vspace{2pt}
\noindent\footnotesize\textsuperscript{a}RRF is undefined when all
payloads pass schema validation on first attempt (SCR\,=\,100\%;
feedback loop was never triggered).
\end{table}

\begin{table}[H]
\centering
\caption{IWSR Breakdown (LLM-Assisted mode)}
\label{tab:breakdown}
\small
\begin{tabular}{llrr}
\toprule
\textbf{Dim.} & \textbf{Category} & \textbf{N} & \textbf{IWSR} \\
\midrule
\multirow{2}{*}{Language}
  & English & 32 & 93.8\% \\
  & Turkish & 28 & 96.4\% \\
\midrule
\multirow{3}{*}{Difficulty}
  & Easy   & 24 & 100.0\% \\
  & Medium & 20 & 100.0\% \\
  & Hard   & 16 &  81.2\% \\
\midrule
\multirow{5}{*}{Service}
  & QoD Session       & 20 &  90.0\% \\
  & Location Retrieval & 15 & 100.0\% \\
  & QoS Profiles      & 10 & 100.0\% \\
  & QoS Provisioning  & 10 & 100.0\% \\
  & Edge Cases        &  5 &  80.0\% \\
\bottomrule
\end{tabular}
\end{table}

Key findings: (1)~SCR and SVR are perfect in both modes — when the
LLM selects the correct service, it always generates a
structurally and semantically valid payload.
(2)~All 5\% IWSR failures in LLM mode are service-selection errors
(HR\,=\,5\%), not payload errors, confirming the effectiveness of
Phase\,2 schema injection.
(3)~Hard-difficulty failures stem from deliberately ambiguous intents
(e.g., ``Apply critical communications QoS'') where the LLM correctly
identifies the QoS priority but misroutes to permanent provisioning
instead of a temporary QoD session.
(4)~Turkish and English intents yield no statistically significant
IWSR difference, validating bilingual robustness.

% =============================================================================
\section{Persistence and Reliability Engineering}
% =============================================================================

\subsection{Persistent NEF Data}

The NEF stack stores location and polygon records in MongoDB. To prevent
data loss across container restarts, the deployment uses a named Docker
volume at \texttt{/data/db} and an auto-seeding initialisation script
mounted at \texttt{/docker-entrypoint-initdb.d/mongo-init.js}.
The first deployment seeds required data automatically; subsequent restarts
preserve the existing dataset without manual reseeding.

\subsection{Three-Layer Validation}

Execution never depends solely on LLM output. All generated payloads pass
three independent validation tiers before any provider call is made:

\begin{itemize}[leftmargin=1.1em]
  \item \textbf{L1 — Schema (SCR):} Pydantic model validation enforcing
        required fields, types, and structural nesting.
  \item \textbf{L2 — Semantic Policy (SVR):} Eight domain-specific rules
        including E.164 phone format, valid QoS profile identifiers,
        duration bounds (1--7200\,s), IPv4 unicast validation, and
        operation/method consistency.
  \item \textbf{L3 — CAPIF Registry (HR):} Live CAPIF discovery query
        to verify the selected service exists in the registry
        (non-blocking if CAPIF is unreachable).
\end{itemize}

This separation between probabilistic LLM reasoning and deterministic
telecom enforcement is the central architectural property enabling
measurable reliability guarantees.

% =============================================================================
\section{Conclusion}
% =============================================================================

This paper presented a working AI-native telco orchestration platform that
integrates 3GPP CAPIF governance, CAMARA service semantics, NEF-backed
telecom realization, three-layer deterministic validation, a structured
feedback correction loop, and multi-provider LLM-assisted planning into one
deployable control plane.
Empirical evaluation on a 60-entry bilingual ground-truth dataset demonstrates
IWSR\,=\,95\% for LLM-assisted mode (Claude Haiku\,4.5) and IWSR\,=\,100\%
for deterministic mode, with perfect schema and semantic compliance in both
cases.
All LLM failures are service-routing errors on deliberately ambiguous
hard-difficulty intents; no payload schema violations occur.
The six reliability metrics (IWSR, SCR, SVR, HR, SS, RRF) constitute a
reusable evaluation framework for any LLM-driven API orchestration system.
The platform source code, ground-truth dataset, and evaluation scripts
are publicly available at
\url{https://github.com/doadurak/mini-telco-platform}.

Future work includes: evaluation with additional LLM providers and model
sizes to establish cross-model reliability bounds; integration of 3GPP ZSM
closed-loop control for autonomous network management; user studies on
intent clarity and its impact on translation reliability; and extension of
the dataset with real operator intent logs for ecological validity.

% =============================================================================
\begin{thebibliography}{10}

\bibitem{3gpp_capif}
3GPP, ``System Architecture for the Common API Framework for 3GPP Northbound
APIs,'' 3GPP TS~23.222, 2023.

\bibitem{camara2022}
J.~Ordonez-Lucena and C.~Dsouza, ``Pathways Towards Network-as-a-Service:
The CAMARA Project,'' \emph{ACM SIGCOMM Workshop on Network-Application
Integration}, 2022.

\bibitem{molina2022}
A.~Molina~Sanchez et al., ``Offering the 3GPP Common API Framework as
Microservice to Vertical Industries,'' \emph{EuCNC/6G Summit}, IEEE, 2022.

\bibitem{nefsim2022}
D.~Fragkos et al., ``NEFSim: An Open Experimentation Framework Utilizing
3GPP's Exposure Services,'' \emph{EuCNC/6G Summit}, IEEE, 2022.

\bibitem{turkovic2025}
B.~Turkovic et al., ``Dynamic QoS Adaptation via Exposed Service APIs for
Advanced CAM Use Cases,'' \emph{EuCNC/6G Summit}, IEEE, 2025.

\bibitem{maatouk2025}
A.~Maatouk et al., ``Large Language Models for Telecom: Forthcoming Impact
on the Industry,'' \emph{IEEE Communications Magazine}, vol.~63, no.~1,
pp.~62--68, Jan.\ 2025.

\bibitem{telagentbench2025}
S.~Lee et al., ``TelAgentBench: A Multi-faceted Benchmark for Evaluating
LLM-based Agents in Telecommunications,'' \emph{EMNLP 2025 Industry Track}.

\bibitem{flowmind2023}
Z.~Zeng et al., ``FlowMind: Automatic Workflow Generation with LLMs,''
\emph{ACM ICAIF}, 2023.

\bibitem{liu2024}
Y.~Liu et al., ``Are LLMs Good at Structured Outputs?,''
\emph{Information Processing and Management}, vol.~61, Art.~103809, 2024.

\bibitem{rafiq2020}
A.~Rafiq et al., ``Intent-Based End-to-End Network Service Orchestration
System for Multi-Platforms,'' \emph{Sustainability}, vol.~12, no.~7,
p.~2782, 2020.

\bibitem{coronado2022}
E.~Coronado et al., ``Zero Touch Management: A Survey of Network Automation
Solutions for 5G and 6G Networks,''
\emph{IEEE Commun.\ Surveys Tuts.}, 2022.

\bibitem{saraiva2019}
N.~F.~Saraiva de~Sousa et al., ``Network Service Orchestration: A Survey,''
\emph{Computer Communications}, vol.~142--143, pp.~69--94, 2019.

\end{thebibliography}

\end{document}
